\printnomenclature[3.5cm] % Значение ширины столбца с обозначениями стоит подбирать вручную
\begin{table}[htbp]
\centering
\caption{Входные множества, параметры и индексы модели}%
\label{tab:theory_notation_input}%
\begin{tabular}{|p{0.18\linewidth}|p{0.74\linewidth}|}
\hline
\textbf{Обозначение} & \textbf{Смысл} \\
% \hline
% \multicolumn{2}{|l|}{\textit{Множества}} \\
\hline
$B$ & множество времён установки/перехода марок бумаги \\
\hline
$C$ & множество продольно-резательных станков (ПРС), включая БРС \\
\hline
$M$ & множество бумагоделательных машин (БДМ) \\
\hline
$O$ & множество заказов \\ %($o_i = (w_i, q_i, h_i, f_i, p_i, l_i, g_i)$) \\
\hline
$T$ & множество складов полуфабрикатов \\
% \hline
% \multicolumn{2}{|l|}{\textit{Количества элементов множеств}} \\
\hline
$I$ & количество БДМ \\
\hline
$J$ & количество ПРС \\
\hline
$K$ & количество складов \\
\hline
$Q$ & количество заказов \\
\hline
$U$ & количество марок бумаги \\
% \hline
% \multicolumn{2}{|l|}{\textit{Прочие параметры}} \\
\hline
$H$ & дата и время начала запуска производства \\
\hline
$V$ & размер кромки тамбура \\
\hline
\end{tabular}
\end{table}

\begin{table}[htbp]
\centering
\caption{Переменные модели}%
\label{tab:theory_notation_output}%
\begin{tabular}{|p{0.18\linewidth}|p{0.74\linewidth}|}
\hline
\textbf{Обозначение} & \textbf{Смысл} \\
\hline
$\overline{\overline{X}}$ & общий план выпуск бумаги на БДМ \\
\hline
$\overline{X}_i$ & план выпуска бумаги на $i$-й БДМ \\
\hline
$X_{i,j}$ & $j$-й тамбур в плане $i$-й БДМ (ширина, марка, плотность) \\
\hline
$\overline{\overline{Y}}$ & общий план нарезки (раскроев) на ПРС/БРС \\
\hline
$\overline{Y}_i$ & план нарезки (раскроев) на $i$-й ПРС/БРС \\
\hline
$Y_{i,j}$ & $j$-й раскрой на $i$-й ПРС (набор форматов, марка, плотность, слои) \\
\hline
$\overline{\overline{Z}}$ & расписание производственных операций выпуска и нарезки \\
\hline
$\overline{Z}$ & производственная операция \\
\hline
$h$, $\overline{h}$ & времена начала/окончания производственной операции \\
\hline
$\overline{T}$ & содержимое складов тамбуров (на момент производственной операции) \\
\hline
$\overline{\overline{T}}$ & зарезервированное содержимое складов (включая тамбуры в производстве) \\
\hline
$G$ & глобальный план $(\overline{\overline{X}}, \overline{\overline{Y}}, \overline{\overline{Z}})$ \\
\hline
\end{tabular}
\end{table}

\begin{table}[htbp]
\centering
\caption{Элементы целевой функции}%
\label{tab:theory_notation_objective}%
\begin{tabular}{|p{0.18\linewidth}|p{0.74\linewidth}|}
\hline
\textbf{Обозначение} & \textbf{Смысл} \\
\hline
$\mu(\overline{\overline{Y}})$ & векторная целевая функция $\langle \omega, \varphi, \tau \rangle \to \min$ \\
\hline
$\omega(\overline{\overline{Y}})$ & суммарная ширина отходов при раскрое \\
\hline
$\varphi(\overline{\overline{Y}})$ & максимальное число раскроев на одной ПРС \\
\hline
$\tau(\overline{\overline{Y}})$ & общее количество перестановок ножей \\
\hline
\end{tabular}
\end{table}
